% --- Complexity Analysis ---
\newcommand{\bigO}[1]{\mathcal{O}(#1)}
\newcommand{\softO}[1]{\tilde{\mathcal{O}}(#1)}
\newcommand{\omegaO}[1]{\Omega(#1)}
\newcommand{\thetaO}[1]{\Theta(#1)}
\newcommand{\smallo}[1]{o(#1)}

% --- Mathematical Sets ---
\newcommand{\RR}{\mathbb{R}}
\newcommand{\ZZ}{\mathbb{Z}}
\newcommand{\NN}{\mathbb{N}}
\newcommand{\QQ}{\mathbb{Q}}

% --- Algorithmic Notations ---
\newcommand{\state}[1]{\text{state}(#1)}
\newcommand{\dpstate}[1]{\text{dp}[#1]}
\newcommand{\cost}[1]{\text{cost}(#1)}
\newcommand{\dist}[1]{\text{dist}(#1)}

% --- Number Theory ---
\newcommand{\floor}[1]{\left\lfloor #1 \right\rfloor}
\newcommand{\ceil}[1]{\left\lceil #1 \right\rceil}
\newcommand{\divisor}{\mid}
\newcommand{\ndivisor}{\nmid}
\newcommand{\modm}[1]{\pmod{#1}}
\newcommand{\congmod}[3]{#1 \equiv #2 \pmod{#3}}
\newcommand{\gcdop}[2]{\text{gcd}(#1, #2)}
\newcommand{\lcmop}[2]{\text{lcm}(#1, #2)}
\newcommand{\eulerphi}[1]{\phi(#1)}
\newcommand{\mobius}[1]{\mu(#1)}
\newcommand{\nCr}[2]{\binom{#1}{#2}}
\newcommand{\nPr}[2]{P_{#2}^{#1}}

% --- Formatting Shorthands ---
\newcommand{\boldtopic}[1]{\textbf{#1}}
\newcommand{\important}[1]{\textcolor{cp-accent}{\textbf{#1}}}
\newcommand{\inlinecode}[1]{\texttt{#1}}

% --- Problem Tags ---
\newcommand{\tagcp}[1]{\hfill \fbox{\small\texttt{#1}}}

% Helper for practice problem links
\newcommand{\practiceproblemlink}[1]{%
    \ifblank{#1}{%
        \textcolor{cp-gray}{\textit{Link no disponible}}%
    }{%
        \scalebox{0.65}{\qrcode{#1}}%
    }%
}

% --- Practice Problems ---
% Usage: \practiceproblem{Name}{Difficulty}{Judge}{Link}
\newcommand{\practiceproblem}[4]{%
    \item \textbf{#1} \hfill \difficultybadge{#2} \\
    \textit{Judge: #3} \hfill \practiceproblemlink{#4}
}

% --- Graph Theory Helpers ---
\newcommand{\tree}[1]{T_{#1}}
\newcommand{\graph}[1]{G_{#1}}
\newcommand{\nodes}{V}
\newcommand{\edges}{E}


% --- Two-column code input ---
% Uso: \twocolumnminted{archivo}{lenguaje}{última línea columna 1}
\newcommand{\twocolumnminted}[3]{%
  \begingroup
    \newcount\startline
    \startline=#3
    \advance\startline by 1
    \begin{columns}
      \begin{column}{0.48\textwidth}
        \inputminted[firstline=1,lastline=#3,linenos,fontsize=\tiny,numbersep=3pt,breaklines]{#2}{#1}
      \end{column}
      \begin{column}{0.48\textwidth}
        \inputminted[firstline=\the\startline,linenos,fontsize=\tiny,numbersep=3pt,breaklines]{#2}{#1}
      \end{column}
    \end{columns}
  \endgroup
}
